\documentclass{article}
\usepackage{pathfinder-readable}

\title{What Fixed-Batch Tail Risk Can Say About Aero Hand Open}

\begin{document}
\maketitle

\begin{abstract}
This account considers Aero Hand Open, a simulated and physical tendon-driven hand, alongside a
paper on fixed-batch tail-risk estimates in a resource-allocation mechanism. The peers asked whether
the second paper could support a claim about the first paper's controller on physical hardware. They
checked a small-sample example and an episode-count bound, but found no recorded physical cube
outcomes with which to assess the controller's tail risk. The thread paused because the mathematical
warning was sound while the proposed link to the hand remained unmeasured.
\end{abstract}

\section{Introduction}

Aero Hand Open presents a cable-driven hand, a simulator that represents its cable transmission, and
a map between simulated and physical actuator commands. It trains a controller in simulation to
rotate a cube and shows rotation on the physical hand without further training. Its training reward
favours rotation and penalises dropping the cube. The physical cube pose is not logged, so the paper
does not report a measured distribution of physical rotation outcomes \cite{Q}.

The other paper studies conditional value at risk, or CVaR, in a quadratic resource-allocation
mechanism. CVaR measures the mean loss in a chosen upper tail. With a fixed number of samples, the
paper's empirical CVaR can favour an allocation whose true CVaR makes participation unattractive,
even when the losses are strongly convex. It also gives conditional bounds for an exact equilibrium
when the expected empirical value stays uniformly close to the true value. Those bounds concern the
mechanism and valid queries within one loss distribution; they do not cover a change from simulation
to hardware \cite{P}.

The peers pursued a narrower connection: could fixed-batch CVaR reasoning support a claim about rare
poor outcomes for the hand's cube-rotation controller? The mechanism paper uses its own labels ``Q''
and ``P'' for other works. Its agents, taxes and opt-out utilities therefore cannot be identified
with parts of the supplied hand paper.

\section{What was done}

The peers first checked the scope of the two papers and set aside a direct transfer of the
mechanism's participation result. They then examined whether a tail-risk estimate could instead
describe complete cube-rotation episodes. One peer proposed comparing physical and simulated episode
losses through their difference. That proposal needed paired, closed-loop outcomes from the cube,
which the hand paper did not provide.

Next, the peers checked prior work. Robot-policy CVaR bounds, including distribution shift and
comparisons among policies, were already studied by Vincent, Feldman and Schwager \cite{vincent}.
Paired real and simulated evaluations for mean performance were also studied in SureSim
\cite{suresim}. These checks removed a generic certificate or a paired comparison, on its own, as a
new result of this thread.

The peers then tested a simpler warning from the mechanism paper: a small batch can make a tail
score misleading. They used an episode that either ends in a fall or does not, and checked both the
CVaR calculation and the number of independent hardware episodes needed for a narrow zero-fall
claim. Finally, they returned to the hand's measured transmission differences and specified the
missing experiment that could connect those differences to cube outcomes.

\section{Results}

Let \(F\) be an episode loss equal to \(1\) for a fall and \(0\) otherwise. Let \(p\) be its fall
probability and \(\beta\) the fraction of the worst outcomes whose mean defines upper-tail CVaR. The
peers derived and cross-checked
\[
\operatorname{CVaR}_{\beta}(F)=\min\{1,p/\beta\}.
\]
With just one sampled episode, empirical CVaR is that episode's loss, so its expected value is
\(p\). At \(p=0.02\) and \(\beta=0.05\), this expected one-episode score is \(0.02\), while true
CVaR is \(0.4\). A second, hypothetical controller with constant loss \(0.1\) looks worse by the
one-episode expected score but better by true CVaR. This is a controller-selection example, not a
measured failure of Aero Hand Open. Its controller is trained on a rotation reward, not selected by
empirical CVaR.

The peers also checked what zero observed falls would say. For \(m\) independent, identically
distributed complete episodes with no falls, the exact one-sided \(95\%\) upper confidence limit on
\(p\) is
\[
p_{\mathrm{upper}}=1-0.05^{1/m}.
\]
To make this limit at most \(0.005\), which would put the corresponding \(5\%\)-tail fall CVaR at
most \(0.1\), requires \(m\geq598\). This calculation assumes a fixed fall definition and comparable
independent episodes. The hand paper gives neither such a hardware episode count nor logged physical
cube pose from which to measure rotation loss.

The hand paper does provide reasons to test the connection. Its two thumb cable excursions differ
from hardware by \(32.7\%\) and \(31.2\%\); simulated actuator responses lead the physical responses
by \(10\) to \(40\) milliseconds. Its attempted randomisation of cube friction and mass is
overwritten. The peers checked that its shared-command replay measures actuator response with the
controller out of the loop. Those measurements identify possible sources of a gap, but they do not
measure the distribution of closed-loop cube rotations or falls \cite{Q}.

\section{What did not hold}

The proposed direct use of the mechanism's result failed on its assumptions: the hand has no
demonstrated counterpart to strategic participation, taxes or opt-out utility. The finite-batch
example establishes a possible evaluation error only. The mechanism paper's value-error bound
concerns an expected empirical CVaR surrogate under valid same-distribution queries. It neither
certifies a particular sampled batch nor bridges simulated and physical losses \cite{P}.

An open-loop actuator replay also cannot settle the task-level question. Similar motor traces can
accompany different object contacts, and the physical cube pose was not logged. The hand's
deployment retunes finger action scales and its nominal pose, so a comparison would have to hold
those settings fixed. The peers distinguished changing the deployed thumb map while keeping the
controller fixed from correcting the simulator and training a new controller. The record contains no
comparison of either intervention on physical episode losses.

\section{Why the thread stopped}

The verifier accepted the Bernoulli CVaR calculation and the \(m\geq598\) bound, but paused the
thread after one round. They illustrate a real measurement issue; they do not show that this
controller has poor tail performance or that its thumb mismatch causes it. Neither paper supplies
the repeated, closed-loop physical loss data needed for that link.

The smallest restart is a specified controller comparison with a declared episode loss and tail
fraction, logged physical cube pose and falls, and independent complete episodes under matched
starting conditions. Keeping the deployment settings fixed would let a thumb-map correction be
compared at a fixed controller. A simulator correction followed by retraining would be a separate
test. Until such outcomes are measured, the pair supports a question and a study design, not a new
hardware tail-risk result.

\bibliographystyle{plainurl}
\bibliography{references}

\end{document}
